% Part: history
% Chapter: biographies
% Section: georg-cantor

\documentclass[../../../include/open-logic-section]{subfiles}

\begin{document}

\olfileid{his}{bio}{can}

\olsection{غيورغ كانتور}

\olphoto{cantor-georg}{Georg Cantor}

روت ترجمة مبكرة لغيورغ كانتور (\textsc{gay}-org
\textsc{kahn}-tor) أنه وُلد في سفينة سائرة إلى سانت بطرسبرغ بروسيا،
ثم عُثر عليه فيها ولم يُعرف أبواه. وهذه رواية لا أصل لها؛ أما الصحيح
أنه وُلد في سانت بطرسبرغ سنة \OLClassicalDigits{1845}.

وحصل كانتور سنة \OLClassicalDigits{1867} على الدكتوراه في الرياضيات من جامعة برلين. واشتهر
بأعماله في نظرية المجموعات حتى نُسب إليه جعلها علمًا مستقلًا للبحث. وهو
أول من أقام البرهان على أن للمجموعات اللانهائية أحجامًا مختلفة. وقد
اشتد الخلاف بين رياضيي زمانه في آرائه، ولا سيما رأيه في مراتب اللانهاية.

وكان اعتقاده الديني ممتزجًا بنظره الرياضي؛ بل زعم أن الله ألقى إليه
مباشرة نظرية الأعداد المتناهية التجاوز. ثم أصابه في أواخر عمره مرض
نفسي، فكان يُدخل المستشفى منذ سنة \OLClassicalDigits{1894}، وكثر ذلك في سنيه الأخيرة.
وساقه النقد العنيف لأعماله، ومنه خصومته للرياضي ليوبولد كرونيكر، إلى
الاكتئاب والعزوف عن الرياضيات. فإذا اشتدت نوبته انصرف إلى الفلسفة
والأدب، ونشر حتى قولًا بأن فرانسيس بيكون هو الذي ألّف مسرحيات شكسبير.

ومات كانتور في \OLClassicalDigits{6} يناير \OLClassicalDigits{1918} بمصح في مدينة هاله.

\begin{reading}
للوقوف على ترجمتين وافيتين لكانتور، انظر \citet{Dauben1990} و
\citet{Grattan-Guinness1971}. وتُعرض آراؤه الجذرية كذلك في برنامج
إذاعة بي بي سي \OLClassicalDigits{4}، \emph{تاريخ موجز للرياضيات}، \citep{Sautoy2014}.
ومن أراد سماع نظرياته في قالب راب فلينظر \citet{Rose2012}.
\end{reading}

\end{document}
